\documentclass[../main.tex]{subfiles}
\begin{document}

\chapter{Datacenter Technologies}
\lessoninfo{
  video={P4L4},
  textbook={Silberschatz et al.\ \cite{silberschatz2018} ch.~1; OSTEP \cite{ostep} ch.~48; Armbrust et al.\ \cite{armbrust2009}; Mell and Grance \cite{mell2011}},
  keywords={Internet service, multi-tier architecture, scale-out, load balancer, homogeneous architecture, heterogeneous architecture, cloud computing, elasticity, law of large numbers, utility computing, IaaS, PaaS, SaaS, multi-tenancy},
}

The preceding lessons presented mechanisms with which processes on
different machines cooperate: remote procedure calls, distributed file
systems and, in Lesson~15, distributed shared memory.  Together with the
threads, servers, scheduling and virtualization of the earlier lessons,
these mechanisms are deployed on the largest scale in datacenters, which run
the services that are used over the Internet.  This final lesson gives a
brief, high-level overview of datacenter technologies: how Internet services
are structured and scaled out over many machines, what cloud computing is,
why it works and which models it offers, what it requires of the underlying
systems, and which technologies make it possible.

\section{Internet services}
\label{sec:l16-services}

Many applications in everyday use do not run on the machines of their
users.\source{P4L4, lesson preview, Internet services.}  They run on
machines in a \term{datacenter}[データセンタ], a facility that houses a
large number of networked computers together with their storage, power
supply and cooling, and operates them to provide computing services.  The
users reach these applications over the Internet.

\begin{definition}[Internet service]
  An \term{Internet service}[インターネットサービス] is any service that is
  provided to its users over the Internet through a web interface.
\end{definition}

Web mail, online shops, search engines and social networks are Internet
services.  Most Internet services consist of three components
(\cref{fig:l16-tiers}):
\begin{description}
  \item[Presentation] interacts with the users.  It receives their requests,
    returns web pages, and provides the \emph{static content}, such as page
    layouts and images, which is the same for every request.
  \item[Business logic] implements the functionality of the service.  It
    processes each request and generates the \emph{dynamic content}, such as
    the results of a search or the contents of a shopping cart.
  \item[Database] stores the data of the service, such as its users,
    products and orders, and retrieves them for the business logic.
\end{description}

\begin{definition}[Multi-tier architecture]
  A service that is structured as a sequence of components, each of which
  uses the services of the next, has a
  \term{multi-tier architecture}[多層アーキテクチャ].  Each component is a
  \emph{tier}; the three components above form a three-tier architecture.
\end{definition}

The tiers are logical components.  They need not be separate processes, and
separate processes need not run on separate machines.  Many open-source and
proprietary middleware products---web servers, application servers,
database servers---implement individual tiers and the interaction between
them, and they can be combined in different configurations.  When the tiers
run as separate processes, the processes communicate through some form of
IPC: message passing or shared memory on the same machine (Lesson~9), and
RPC or RMI across machines (Lesson~13).

\begin{figure}[htbp]
  \centering
  % Three-tier Internet service: (a) all tiers on one machine, (b) each tier
% on its own machine, the tiers communicating by RPC/RMI.
% Usage: % Three-tier Internet service: (a) all tiers on one machine, (b) each tier
% on its own machine, the tiers communicating by RPC/RMI.
% Usage: % Three-tier Internet service: (a) all tiers on one machine, (b) each tier
% on its own machine, the tiers communicating by RPC/RMI.
% Usage: \input{figures/l16-tiers}   (width ~118 mm)
\begin{tikzpicture}[x=1mm, y=1mm, font=\scriptsize\sffamily,
  >={Stealth[length=1.5mm]},
  tier/.style={draw, fill=ln-box, minimum width=22mm, minimum height=10mm, inner sep=1pt, align=center},
  cl/.style={draw, fill=white, minimum width=14mm, minimum height=10mm, inner sep=1pt, align=center},
  machine/.style={draw, rounded corners=0.8mm},
  lab/.style={font=\scriptsize\sffamily, text=ln-muted, align=center},
  who/.style={font=\scriptsize\sffamily, text=ln-muted, anchor=north west, inner sep=0.8pt},
  ttl/.style={font=\footnotesize\sffamily\bfseries, anchor=west}]
  % ---------------- (a) one machine
  \node[ttl] at (0,57) {\iflangja{(a) すべての層が1台のマシン上にある}{(a) all tiers on one machine}};
  \node[cl] (c1) at (7,40) {\iflangja{クライアント}{clients}};
  \draw[machine] (20,30) rectangle (118,53);
  \node[who] at (20.5,52.5) {\iflangja{マシン}{machine}};
  \node[tier] (p1) at (33,42) {\iflangja{プレゼンテーション}{presentation}\\[0.3mm]\textcolor{ln-muted}{\iflangja{静的コンテンツ}{static content}}};
  \node[tier] (b1) at (69,42) {\iflangja{ビジネスロジック}{business logic}\\[0.3mm]\textcolor{ln-muted}{\iflangja{動的コンテンツ}{dynamic content}}};
  \node[tier] (d1) at (105,42) {\iflangja{データベース}{database}\\[0.3mm]\textcolor{ln-muted}{\iflangja{データストア}{data store}}};
  \draw[<->, thick] (c1.east |- p1) -- (p1.west);
  \draw[<->, thick] (p1) -- (b1);
  \draw[<->, thick] (b1) -- (d1);
  \node[lab] at (69,33.2)
    {\iflangja{1つのプロセス，または IPC・共有メモリで通信する複数のプロセス}{one process, or several processes communicating by IPC or shared memory}};
  % ---------------- (b) separate machines
  \node[ttl] at (0,24) {\iflangja{(b) 各層が別のマシン上にある}{(b) each tier on its own machine}};
  \node[cl] (c2) at (7,9) {\iflangja{クライアント}{clients}};
  \foreach \x/\n in {33/{\iflangja{Web サーバ}{web server}}, 69/{\iflangja{アプリケーションサーバ}{application server}}, 105/{\iflangja{データベースサーバ}{database server}}} {
    % Japanese: 2 mm wider boxes, or "アプリケーションサーバ" touches the frame.
    \draw[machine] (\x-\iflangja{14}{13},0) rectangle (\x+\iflangja{14}{13},20);
    \node[who] at (\x-12.5,19.5) {\n};
  }
  \node[tier] (p2) at (33,8) {\iflangja{プレゼンテーション}{presentation}};
  \node[tier] (b2) at (69,8) {\iflangja{ビジネスロジック}{business logic}};
  \node[tier] (d2) at (105,8) {\iflangja{データベース}{database}};
  \draw[<->, thick] (c2.east |- p2) -- (p2.west);
  \draw[<->, thick] (p2) -- (b2);
  \draw[<->, thick] (b2) -- (d2);
  \node[lab, anchor=south] at (51,8.6) {RPC};
  \node[lab, anchor=north] at (51,7.4) {RMI};
  \node[lab, anchor=south] at (87,8.6) {RPC};
  \node[lab, anchor=north] at (87,7.4) {RMI};
\end{tikzpicture}
   (width ~118 mm)
\begin{tikzpicture}[x=1mm, y=1mm, font=\scriptsize\sffamily,
  >={Stealth[length=1.5mm]},
  tier/.style={draw, fill=ln-box, minimum width=22mm, minimum height=10mm, inner sep=1pt, align=center},
  cl/.style={draw, fill=white, minimum width=14mm, minimum height=10mm, inner sep=1pt, align=center},
  machine/.style={draw, rounded corners=0.8mm},
  lab/.style={font=\scriptsize\sffamily, text=ln-muted, align=center},
  who/.style={font=\scriptsize\sffamily, text=ln-muted, anchor=north west, inner sep=0.8pt},
  ttl/.style={font=\footnotesize\sffamily\bfseries, anchor=west}]
  % ---------------- (a) one machine
  \node[ttl] at (0,57) {\iflangja{(a) すべての層が1台のマシン上にある}{(a) all tiers on one machine}};
  \node[cl] (c1) at (7,40) {\iflangja{クライアント}{clients}};
  \draw[machine] (20,30) rectangle (118,53);
  \node[who] at (20.5,52.5) {\iflangja{マシン}{machine}};
  \node[tier] (p1) at (33,42) {\iflangja{プレゼンテーション}{presentation}\\[0.3mm]\textcolor{ln-muted}{\iflangja{静的コンテンツ}{static content}}};
  \node[tier] (b1) at (69,42) {\iflangja{ビジネスロジック}{business logic}\\[0.3mm]\textcolor{ln-muted}{\iflangja{動的コンテンツ}{dynamic content}}};
  \node[tier] (d1) at (105,42) {\iflangja{データベース}{database}\\[0.3mm]\textcolor{ln-muted}{\iflangja{データストア}{data store}}};
  \draw[<->, thick] (c1.east |- p1) -- (p1.west);
  \draw[<->, thick] (p1) -- (b1);
  \draw[<->, thick] (b1) -- (d1);
  \node[lab] at (69,33.2)
    {\iflangja{1つのプロセス，または IPC・共有メモリで通信する複数のプロセス}{one process, or several processes communicating by IPC or shared memory}};
  % ---------------- (b) separate machines
  \node[ttl] at (0,24) {\iflangja{(b) 各層が別のマシン上にある}{(b) each tier on its own machine}};
  \node[cl] (c2) at (7,9) {\iflangja{クライアント}{clients}};
  \foreach \x/\n in {33/{\iflangja{Web サーバ}{web server}}, 69/{\iflangja{アプリケーションサーバ}{application server}}, 105/{\iflangja{データベースサーバ}{database server}}} {
    % Japanese: 2 mm wider boxes, or "アプリケーションサーバ" touches the frame.
    \draw[machine] (\x-\iflangja{14}{13},0) rectangle (\x+\iflangja{14}{13},20);
    \node[who] at (\x-12.5,19.5) {\n};
  }
  \node[tier] (p2) at (33,8) {\iflangja{プレゼンテーション}{presentation}};
  \node[tier] (b2) at (69,8) {\iflangja{ビジネスロジック}{business logic}};
  \node[tier] (d2) at (105,8) {\iflangja{データベース}{database}};
  \draw[<->, thick] (c2.east |- p2) -- (p2.west);
  \draw[<->, thick] (p2) -- (b2);
  \draw[<->, thick] (b2) -- (d2);
  \node[lab, anchor=south] at (51,8.6) {RPC};
  \node[lab, anchor=north] at (51,7.4) {RMI};
  \node[lab, anchor=south] at (87,8.6) {RPC};
  \node[lab, anchor=north] at (87,7.4) {RMI};
\end{tikzpicture}
   (width ~118 mm)
\begin{tikzpicture}[x=1mm, y=1mm, font=\scriptsize\sffamily,
  >={Stealth[length=1.5mm]},
  tier/.style={draw, fill=ln-box, minimum width=22mm, minimum height=10mm, inner sep=1pt, align=center},
  cl/.style={draw, fill=white, minimum width=14mm, minimum height=10mm, inner sep=1pt, align=center},
  machine/.style={draw, rounded corners=0.8mm},
  lab/.style={font=\scriptsize\sffamily, text=ln-muted, align=center},
  who/.style={font=\scriptsize\sffamily, text=ln-muted, anchor=north west, inner sep=0.8pt},
  ttl/.style={font=\footnotesize\sffamily\bfseries, anchor=west}]
  % ---------------- (a) one machine
  \node[ttl] at (0,57) {\iflangja{(a) すべての層が1台のマシン上にある}{(a) all tiers on one machine}};
  \node[cl] (c1) at (7,40) {\iflangja{クライアント}{clients}};
  \draw[machine] (20,30) rectangle (118,53);
  \node[who] at (20.5,52.5) {\iflangja{マシン}{machine}};
  \node[tier] (p1) at (33,42) {\iflangja{プレゼンテーション}{presentation}\\[0.3mm]\textcolor{ln-muted}{\iflangja{静的コンテンツ}{static content}}};
  \node[tier] (b1) at (69,42) {\iflangja{ビジネスロジック}{business logic}\\[0.3mm]\textcolor{ln-muted}{\iflangja{動的コンテンツ}{dynamic content}}};
  \node[tier] (d1) at (105,42) {\iflangja{データベース}{database}\\[0.3mm]\textcolor{ln-muted}{\iflangja{データストア}{data store}}};
  \draw[<->, thick] (c1.east |- p1) -- (p1.west);
  \draw[<->, thick] (p1) -- (b1);
  \draw[<->, thick] (b1) -- (d1);
  \node[lab] at (69,33.2)
    {\iflangja{1つのプロセス，または IPC・共有メモリで通信する複数のプロセス}{one process, or several processes communicating by IPC or shared memory}};
  % ---------------- (b) separate machines
  \node[ttl] at (0,24) {\iflangja{(b) 各層が別のマシン上にある}{(b) each tier on its own machine}};
  \node[cl] (c2) at (7,9) {\iflangja{クライアント}{clients}};
  \foreach \x/\n in {33/{\iflangja{Web サーバ}{web server}}, 69/{\iflangja{アプリケーションサーバ}{application server}}, 105/{\iflangja{データベースサーバ}{database server}}} {
    % Japanese: 2 mm wider boxes, or "アプリケーションサーバ" touches the frame.
    \draw[machine] (\x-\iflangja{14}{13},0) rectangle (\x+\iflangja{14}{13},20);
    \node[who] at (\x-12.5,19.5) {\n};
  }
  \node[tier] (p2) at (33,8) {\iflangja{プレゼンテーション}{presentation}};
  \node[tier] (b2) at (69,8) {\iflangja{ビジネスロジック}{business logic}};
  \node[tier] (d2) at (105,8) {\iflangja{データベース}{database}};
  \draw[<->, thick] (c2.east |- p2) -- (p2.west);
  \draw[<->, thick] (p2) -- (b2);
  \draw[<->, thick] (b2) -- (d2);
  \node[lab, anchor=south] at (51,8.6) {RPC};
  \node[lab, anchor=north] at (51,7.4) {RMI};
  \node[lab, anchor=south] at (87,8.6) {RPC};
  \node[lab, anchor=north] at (87,7.4) {RMI};
\end{tikzpicture}

  \caption{A three-tier Internet service.  (a) All tiers run on one machine,
    in one process or in several processes.  (b) Every tier runs on its own
    machine, and the tiers communicate by RPC or RMI.}
  \label{fig:l16-tiers}
\end{figure}

\section{Scale-out architectures}
\label{sec:l16-scaleout}

A single machine serves only a limited number of requests per second,
however efficiently its server software is designed (Lesson~6).\source{P4L4,
Internet service architectures, homogeneous architectures, heterogeneous
architectures.}  A service that must handle more requests therefore runs as
many processes on many machines, the \emph{nodes} of the service.

\begin{definition}[Scale-out]
  An architecture that increases the capacity of a service by adding nodes,
  rather than by replacing a node with a more powerful machine
  (\emph{scale-up}), is a \term{scale-out}[スケールアウト] architecture.
\end{definition}

Scale-out architectures commonly apply the boss-workers
pattern\index{boss-workers pattern} of Lesson~3 to machines instead of
threads: a front-end node receives all requests and passes each of them to
one of the nodes that process requests.

\begin{definition}[Load balancer]
  A \term{load balancer}[ロードバランサ] is a front end that distributes the
  incoming requests of a service among the nodes that process them.
\end{definition}

Just as the workers of the boss-workers pattern are either all equal or
specialized, the nodes of a scale-out architecture can be organized in two
ways (\cref{fig:l16-scaleout}):
\begin{description}
  \item[All equal] Every node can execute every step of the processing of
    every request.  The nodes are functionally homogeneous.
  \item[Specialized] Each node executes only some steps of the processing,
    or only requests of some types.  The nodes are functionally
    heterogeneous.
\end{description}

\begin{figure}[htbp]
  \centering
  % Scale-out architectures with a front-end load balancer: (a) functionally
% homogeneous nodes, (b) functionally heterogeneous (specialized) nodes.
% Usage: % Scale-out architectures with a front-end load balancer: (a) functionally
% homogeneous nodes, (b) functionally heterogeneous (specialized) nodes.
% Usage: % Scale-out architectures with a front-end load balancer: (a) functionally
% homogeneous nodes, (b) functionally heterogeneous (specialized) nodes.
% Usage: \input{figures/l16-scaleout}   (width ~118 mm)
\begin{tikzpicture}[x=1mm, y=1mm, font=\scriptsize\sffamily,
  >={Stealth[length=1.5mm]},
  fe/.style={draw, fill=ln-accent!18, minimum width=30mm, minimum height=8mm, inner sep=1pt, align=center},
  nd/.style={draw, fill=ln-box, minimum width=12mm, minimum height=9mm, inner sep=0.5pt, align=center},
  data/.style={draw, fill=ln-accent!35, minimum height=6mm, inner sep=1pt, align=center},
  lab/.style={font=\scriptsize\sffamily, text=ln-muted, align=center},
  ttl/.style={font=\footnotesize\sffamily\bfseries, anchor=west}]
  % ---------------- (a) homogeneous
  \begin{scope}[shift={(0,0)}]
    \node[ttl] at (0,56) {\iflangja{(a) ホモジニアス構成}{(a) homogeneous}};
    \draw[->, thick] (28.5,51) -- node[lab, right] {\iflangja{要求}{requests}} (28.5,44.5);
    \node[fe] (f) at (28.5,40) {\iflangja{フロントエンド}{front end}\\\textcolor{ln-muted}{\iflangja{要求を順に割り当てる}{assigns in turn}}};
    \foreach \i/\x in {1/7.5,2/21.5,3/35.5,4/49.5} {
      \node[nd] (n\i) at (\x,22) {\iflangja{ノード}{node}\\\textcolor{ln-muted}{A, B, C}};
      \draw[->] (f.south) -- (n\i.north);
      \draw (n\i.south) -- (\x,10);
    }
    \node[data, minimum width=54mm] at (28.5,7) {\iflangja{データ（全ノードからアクセス可能）}{data (accessible from every node)}};
  \end{scope}
  % ---------------- (b) heterogeneous
  \begin{scope}[shift={(61,0)}]
    \node[ttl] at (0,56) {\iflangja{(b) ヘテロジニアス構成}{(b) heterogeneous}};
    \draw[->, thick] (28.5,51) -- node[lab, right] {\iflangja{要求}{requests}} (28.5,44.5);
    \node[fe] (f) at (28.5,40) {\iflangja{フロントエンド}{front end}\\\textcolor{ln-muted}{\iflangja{要求の種類を調べる}{inspects each request}}};
    \foreach \i/\x/\t in {1/7.5/A,2/21.5/A,3/35.5/B,4/49.5/C} {
      \node[nd] (n\i) at (\x,22) {\iflangja{ノード}{node}\\\textcolor{ln-muted}{\t}};
      \draw[->] (f.south) -- (n\i.north);
      \draw (n\i.south) -- (\x,10);
    }
    \node[data, minimum width=26mm] at (14.5,7) {\iflangja{データ}{data} A};
    \node[data, minimum width=12mm] at (35.5,7) {B};
    \node[data, minimum width=12mm] at (49.5,7) {C};
  \end{scope}
  \node[lab, anchor=north] at (59,1.5)
    {\iflangja{A，B，C：要求の種類}{A, B, C: request types}};
\end{tikzpicture}
   (width ~118 mm)
\begin{tikzpicture}[x=1mm, y=1mm, font=\scriptsize\sffamily,
  >={Stealth[length=1.5mm]},
  fe/.style={draw, fill=ln-accent!18, minimum width=30mm, minimum height=8mm, inner sep=1pt, align=center},
  nd/.style={draw, fill=ln-box, minimum width=12mm, minimum height=9mm, inner sep=0.5pt, align=center},
  data/.style={draw, fill=ln-accent!35, minimum height=6mm, inner sep=1pt, align=center},
  lab/.style={font=\scriptsize\sffamily, text=ln-muted, align=center},
  ttl/.style={font=\footnotesize\sffamily\bfseries, anchor=west}]
  % ---------------- (a) homogeneous
  \begin{scope}[shift={(0,0)}]
    \node[ttl] at (0,56) {\iflangja{(a) ホモジニアス構成}{(a) homogeneous}};
    \draw[->, thick] (28.5,51) -- node[lab, right] {\iflangja{要求}{requests}} (28.5,44.5);
    \node[fe] (f) at (28.5,40) {\iflangja{フロントエンド}{front end}\\\textcolor{ln-muted}{\iflangja{要求を順に割り当てる}{assigns in turn}}};
    \foreach \i/\x in {1/7.5,2/21.5,3/35.5,4/49.5} {
      \node[nd] (n\i) at (\x,22) {\iflangja{ノード}{node}\\\textcolor{ln-muted}{A, B, C}};
      \draw[->] (f.south) -- (n\i.north);
      \draw (n\i.south) -- (\x,10);
    }
    \node[data, minimum width=54mm] at (28.5,7) {\iflangja{データ（全ノードからアクセス可能）}{data (accessible from every node)}};
  \end{scope}
  % ---------------- (b) heterogeneous
  \begin{scope}[shift={(61,0)}]
    \node[ttl] at (0,56) {\iflangja{(b) ヘテロジニアス構成}{(b) heterogeneous}};
    \draw[->, thick] (28.5,51) -- node[lab, right] {\iflangja{要求}{requests}} (28.5,44.5);
    \node[fe] (f) at (28.5,40) {\iflangja{フロントエンド}{front end}\\\textcolor{ln-muted}{\iflangja{要求の種類を調べる}{inspects each request}}};
    \foreach \i/\x/\t in {1/7.5/A,2/21.5/A,3/35.5/B,4/49.5/C} {
      \node[nd] (n\i) at (\x,22) {\iflangja{ノード}{node}\\\textcolor{ln-muted}{\t}};
      \draw[->] (f.south) -- (n\i.north);
      \draw (n\i.south) -- (\x,10);
    }
    \node[data, minimum width=26mm] at (14.5,7) {\iflangja{データ}{data} A};
    \node[data, minimum width=12mm] at (35.5,7) {B};
    \node[data, minimum width=12mm] at (49.5,7) {C};
  \end{scope}
  \node[lab, anchor=north] at (59,1.5)
    {\iflangja{A，B，C：要求の種類}{A, B, C: request types}};
\end{tikzpicture}
   (width ~118 mm)
\begin{tikzpicture}[x=1mm, y=1mm, font=\scriptsize\sffamily,
  >={Stealth[length=1.5mm]},
  fe/.style={draw, fill=ln-accent!18, minimum width=30mm, minimum height=8mm, inner sep=1pt, align=center},
  nd/.style={draw, fill=ln-box, minimum width=12mm, minimum height=9mm, inner sep=0.5pt, align=center},
  data/.style={draw, fill=ln-accent!35, minimum height=6mm, inner sep=1pt, align=center},
  lab/.style={font=\scriptsize\sffamily, text=ln-muted, align=center},
  ttl/.style={font=\footnotesize\sffamily\bfseries, anchor=west}]
  % ---------------- (a) homogeneous
  \begin{scope}[shift={(0,0)}]
    \node[ttl] at (0,56) {\iflangja{(a) ホモジニアス構成}{(a) homogeneous}};
    \draw[->, thick] (28.5,51) -- node[lab, right] {\iflangja{要求}{requests}} (28.5,44.5);
    \node[fe] (f) at (28.5,40) {\iflangja{フロントエンド}{front end}\\\textcolor{ln-muted}{\iflangja{要求を順に割り当てる}{assigns in turn}}};
    \foreach \i/\x in {1/7.5,2/21.5,3/35.5,4/49.5} {
      \node[nd] (n\i) at (\x,22) {\iflangja{ノード}{node}\\\textcolor{ln-muted}{A, B, C}};
      \draw[->] (f.south) -- (n\i.north);
      \draw (n\i.south) -- (\x,10);
    }
    \node[data, minimum width=54mm] at (28.5,7) {\iflangja{データ（全ノードからアクセス可能）}{data (accessible from every node)}};
  \end{scope}
  % ---------------- (b) heterogeneous
  \begin{scope}[shift={(61,0)}]
    \node[ttl] at (0,56) {\iflangja{(b) ヘテロジニアス構成}{(b) heterogeneous}};
    \draw[->, thick] (28.5,51) -- node[lab, right] {\iflangja{要求}{requests}} (28.5,44.5);
    \node[fe] (f) at (28.5,40) {\iflangja{フロントエンド}{front end}\\\textcolor{ln-muted}{\iflangja{要求の種類を調べる}{inspects each request}}};
    \foreach \i/\x/\t in {1/7.5/A,2/21.5/A,3/35.5/B,4/49.5/C} {
      \node[nd] (n\i) at (\x,22) {\iflangja{ノード}{node}\\\textcolor{ln-muted}{\t}};
      \draw[->] (f.south) -- (n\i.north);
      \draw (n\i.south) -- (\x,10);
    }
    \node[data, minimum width=26mm] at (14.5,7) {\iflangja{データ}{data} A};
    \node[data, minimum width=12mm] at (35.5,7) {B};
    \node[data, minimum width=12mm] at (49.5,7) {C};
  \end{scope}
  \node[lab, anchor=north] at (59,1.5)
    {\iflangja{A，B，C：要求の種類}{A, B, C: request types}};
\end{tikzpicture}

  \caption{Scale-out architectures.  (a) Every node processes requests of
    every type, and every node can access all data.  (b) Nodes are
    specialized for request types, and the front end must inspect each
    request to choose a node.}
  \label{fig:l16-scaleout}
\end{figure}

\subsection{Homogeneous architectures}

\begin{definition}[Homogeneous architecture]
  In a \term{homogeneous architecture}[ホモジニアス構成] all nodes are
  functionally identical: each of them can perform every processing step
  for every request.
\end{definition}

The front end of a homogeneous architecture is simple.  It need not examine
the requests; it can assign them to the nodes in turn, in round-robin order,
or to the node with the least load.  Functional homogeneity does not require
every node to store all data.  It requires only that every node can
\emph{access} all data, for instance through a distributed file system or a
shared database tier, which may itself be replicated or partitioned
(Lesson~14).  To increase the capacity, nodes are simply added.

The price is that the service benefits little from caching.  The front end
does not know which node has recently processed a similar request and holds
the data it needs in a hot cache\index{hot cache}.  A request that one node
could serve from its cache is therefore likely to be sent to a node that must
first fetch the data---the lack of locality that Lesson~3 noted for equal
workers.

\subsection{Heterogeneous architectures}

\begin{definition}[Heterogeneous architecture]
  In a \term{heterogeneous architecture}[ヘテロジニアス構成] the nodes are
  functionally specialized: different nodes perform different processing
  steps or serve different types of requests.
\end{definition}

Data then need not be uniformly accessible from all nodes; each node needs
only the data of its own processing steps or request types.  The advantage
is locality.  All requests of a given type reach the same few nodes, which
keep the code and data these requests need in hot caches, as the specialized
threads of Lesson~3\index{specialization} do.  The front end, however,
becomes more complex: it must inspect every request, for example parse its
URL, to determine which nodes can process it.  The management of the service
becomes more complex as well.

\subsection{Scaling and management}

The two architectures differ most when the load grows.  To scale a
homogeneous architecture, it suffices to add nodes, together with the
processes, machines and access to storage that they need.  Every new node
takes its share of every type of request, and no knowledge of which requests
cause the load is required.  To scale a heterogeneous architecture, the
operator must know which types of requests are in demand and add nodes of
these types, which requires profiling the workload\index{workload}.  As the
demand for the different request types changes over time, the mix of nodes
must follow it.  Otherwise the nodes of a popular type become \emph{hot
spots} with long queues of pending requests, while the nodes of other types
are idle.  \Cref{tab:l16-arch} summarizes the comparison.

\begin{table}[htbp]
  \centering
  \caption{Homogeneous and heterogeneous scale-out architectures.}
  \label{tab:l16-arch}
  \small
  \begin{tabular}{@{}>{\raggedright\arraybackslash}p{0.2\linewidth}>{\raggedright\arraybackslash}p{0.36\linewidth}>{\raggedright\arraybackslash}p{0.36\linewidth}@{}}
    \toprule
    & Homogeneous & Heterogeneous \\
    \midrule
    Nodes & functionally identical & specialized by processing step or
      request type \\
    Front end & simple; need not inspect requests & must inspect every
      request \\
    Data & accessible from every node & needed only on some nodes \\
    Caching and locality & little benefit & requests of a type reach the
      same nodes \\
    Scaling & add nodes & profile the demand; add nodes of the types
      needed \\
    Change in request mix & absorbed without intervention & hot spots
      unless nodes are reassigned \\
    \bottomrule
  \end{tabular}
\end{table}

\begin{example}
  \label{ex:l16-scaling}
  An online shop receives two types of requests.  On a node of a
  homogeneous architecture, a search takes \SI{5}{\milli\second} of
  processing and a checkout \SI{20}{\milli\second}.  A load of 1200 searches
  and 200 checkouts per second requires $1200 \times 0.005 + 200 \times 0.020
  = 6 + 4 = 10$ seconds of processing per second, the capacity of ten nodes.
  When the mix changes to 600 searches and 350 checkouts per second, the
  requirement is $3 + 7 = 10$ again: the same ten nodes suffice, and nobody
  needs to notice the change.

  In a heterogeneous architecture, suppose that hot caches let the
  specialized nodes process requests in \SI{20}{\percent} less time,
  \SI{4}{\milli\second} per search and \SI{16}{\milli\second} per checkout.
  The original load needs $\lceil 1200 \times 0.004 \rceil = 5$ search nodes
  and $\lceil 200 \times 0.016 \rceil = 4$ checkout nodes, nine in total.
  After the change, $\lceil 2.4 \rceil = 3$ search nodes and
  $\lceil 5.6 \rceil = 6$ checkout nodes are needed---still nine, but in a
  different mix.  Until the operator moves two nodes from search to
  checkout, the four checkout nodes complete only $4/0.016 = 250$ checkouts
  per second, the queue of pending checkouts grows by 100 requests every
  second, and two of the five search nodes are idle.
\end{example}

\begin{keypoint}
  An Internet service is organized in tiers and scaled out over many nodes
  behind a load balancer.  Homogeneous nodes keep the front end and scaling
  simple but gain little from caching; heterogeneous nodes gain locality at
  the cost of a front end that inspects requests and of management that must
  follow the mix of requests.
\end{keypoint}

\section{The motivation for cloud computing}
\label{sec:l16-motivation}

Amazon, the online retailer, had provisioned its hardware for the peak load
of the holiday shopping season, and for the rest of the year much of this
hardware was idle.\source{P4L4, cloud computing poster child: Animoto, cloud
computing requirements; Armbrust et al.\ \cite{armbrust2009}.}  Amazon
opened access to it through web-based APIs, so that third parties could run
their own workloads on Amazon's hardware for a fee.  The result was Amazon
Web Services (AWS).  Its Elastic Compute Cloud (EC2), launched in 2006, rents
out \emph{compute instances}: virtual machines (Lesson~12) that customers
start and stop through the API and pay for by the hour.

Animoto, a service that produces videos from users' photographs, rented its
compute instances from EC2.  In April 2008 Animoto became available to the
users of Facebook, and its demand surged: within three days the number of
servers it used grew from 50 to 3500, after which its traffic fell to well
below the peak \cite{armbrust2009}.  Growth by almost two orders of magnitude
within days is impossible when machines must be bought, installed and
configured; with rented resources it required only calls to an API.

\subsection{Provisioning for variable demand}

Traditionally, a service buys and configures its own resources.  Their
capacity is determined in advance from the expected demand, usually from its
expected peak, and it can later be increased only in large steps, each
delayed by a purchase.  Since demand varies over time, a fixed capacity is
wrong most of the time (\cref{fig:l16-capacity}a):
\begin{itemize}
  \item while the demand is below the capacity, the idle resources are
    wasted: they have been paid for but are not used;
  \item when the demand exceeds the capacity, requests are dropped, and the
    opportunity to serve these users, and the revenue they would have
    brought, is lost.
\end{itemize}

\begin{figure}[htbp]
  \centering
  % Capacity versus demand over time: (a) fixed capacity provisioned for the
% expected peak wastes resources and loses requests during an unexpected
% surge; (b) elastic capacity follows the demand.
% Usage: % Capacity versus demand over time: (a) fixed capacity provisioned for the
% expected peak wastes resources and loses requests during an unexpected
% surge; (b) elastic capacity follows the demand.
% Usage: % Capacity versus demand over time: (a) fixed capacity provisioned for the
% expected peak wastes resources and loses requests during an unexpected
% surge; (b) elastic capacity follows the demand.
% Usage: \input{figures/l16-capacity}   (width ~118 mm)
\begin{tikzpicture}[x=1mm, y=1mm, font=\scriptsize\sffamily,
  >={Stealth[length=1.5mm]},
  ax/.style={->, thin},
  dem/.style={thick, ln-accent},
  capa/.style={thick, ln-caution},
  lab/.style={font=\scriptsize\sffamily, align=center},
  ttl/.style={font=\footnotesize\sffamily\bfseries, anchor=west}]
  % ---------------- (a) fixed capacity
  \begin{scope}[shift={(3,0)}]
    \node[ttl] at (-4,43) {\iflangja{(a) 固定された容量}{(a) fixed capacity}};
    \begin{scope}
      \clip (0,0) rectangle (50,22);
      \fill[ln-rule!45] (0,34) -- plot[smooth, tension=0.6] coordinates
        {(0,6) (5,8) (10,12) (15,10) (20,14) (25,17) (30,15) (35,19) (38,29) (41,28) (44,18) (50,16)}
        -- (50,34) -- cycle;
    \end{scope}
    \begin{scope}
      \clip (0,22) rectangle (50,34);
      \fill[ln-caution!30] (0,0) -- plot[smooth, tension=0.6] coordinates
        {(0,6) (5,8) (10,12) (15,10) (20,14) (25,17) (30,15) (35,19) (38,29) (41,28) (44,18) (50,16)}
        -- (50,0) -- cycle;
    \end{scope}
    \draw[dem] plot[smooth, tension=0.6] coordinates
      {(0,6) (5,8) (10,12) (15,10) (20,14) (25,17) (30,15) (35,19) (38,29) (41,28) (44,18) (50,16)};
    \draw[capa] (0,22) -- (50,22);
    \draw[ax] (0,0) -- (53,0) node[below left, lab] {\iflangja{時間}{time}};
    \draw[ax] (0,0) -- (0,35) node[above right, lab, anchor=south west] {\iflangja{資源}{resources}};
    \node[lab, fill=white, inner sep=0.6pt] at (14,18.5) {\iflangja{無駄な資源}{wasted resources}};
    \draw[->, ln-caution] (30,30) -- (37,25.5);
    \node[lab, anchor=east, text=ln-caution] at (30.5,30.5) {\iflangja{逸失した機会}{lost opportunity}};
  \end{scope}
  % ---------------- (b) elastic capacity
  \begin{scope}[shift={(63,0)}]
    \node[ttl] at (-4,43) {\iflangja{(b) 伸縮する容量}{(b) elastic capacity}};
    \draw[dem] plot[smooth, tension=0.6] coordinates
      {(0,6) (5,8) (10,12) (15,10) (20,14) (25,17) (30,15) (35,19) (38,29) (41,28) (44,18) (50,16)};
    \draw[capa] (0,9) -- (5,9) -- (5,13.5) -- (10,13.5) -- (10,14) -- (15,14)
      -- (15,15.5) -- (20,15.5) -- (20,19) -- (30,19) -- (30,21) -- (35,21)
      -- (35,31) -- (43,31) -- (43,21) -- (46,21) -- (46,18.5) -- (50,18.5);
    \draw[ax] (0,0) -- (53,0) node[below left, lab] {\iflangja{時間}{time}};
    \draw[ax] (0,0) -- (0,35) node[above right, lab, anchor=south west] {\iflangja{資源}{resources}};
  \end{scope}
  % ---------------- legend
  \draw[dem] (36,-7) -- (44,-7);
  \node[lab, anchor=west] at (45,-7) {\iflangja{需要}{demand}};
  \draw[capa] (62,-7) -- (70,-7);
  \node[lab, anchor=west] at (71,-7) {\iflangja{容量}{capacity}};
\end{tikzpicture}
   (width ~118 mm)
\begin{tikzpicture}[x=1mm, y=1mm, font=\scriptsize\sffamily,
  >={Stealth[length=1.5mm]},
  ax/.style={->, thin},
  dem/.style={thick, ln-accent},
  capa/.style={thick, ln-caution},
  lab/.style={font=\scriptsize\sffamily, align=center},
  ttl/.style={font=\footnotesize\sffamily\bfseries, anchor=west}]
  % ---------------- (a) fixed capacity
  \begin{scope}[shift={(3,0)}]
    \node[ttl] at (-4,43) {\iflangja{(a) 固定された容量}{(a) fixed capacity}};
    \begin{scope}
      \clip (0,0) rectangle (50,22);
      \fill[ln-rule!45] (0,34) -- plot[smooth, tension=0.6] coordinates
        {(0,6) (5,8) (10,12) (15,10) (20,14) (25,17) (30,15) (35,19) (38,29) (41,28) (44,18) (50,16)}
        -- (50,34) -- cycle;
    \end{scope}
    \begin{scope}
      \clip (0,22) rectangle (50,34);
      \fill[ln-caution!30] (0,0) -- plot[smooth, tension=0.6] coordinates
        {(0,6) (5,8) (10,12) (15,10) (20,14) (25,17) (30,15) (35,19) (38,29) (41,28) (44,18) (50,16)}
        -- (50,0) -- cycle;
    \end{scope}
    \draw[dem] plot[smooth, tension=0.6] coordinates
      {(0,6) (5,8) (10,12) (15,10) (20,14) (25,17) (30,15) (35,19) (38,29) (41,28) (44,18) (50,16)};
    \draw[capa] (0,22) -- (50,22);
    \draw[ax] (0,0) -- (53,0) node[below left, lab] {\iflangja{時間}{time}};
    \draw[ax] (0,0) -- (0,35) node[above right, lab, anchor=south west] {\iflangja{資源}{resources}};
    \node[lab, fill=white, inner sep=0.6pt] at (14,18.5) {\iflangja{無駄な資源}{wasted resources}};
    \draw[->, ln-caution] (30,30) -- (37,25.5);
    \node[lab, anchor=east, text=ln-caution] at (30.5,30.5) {\iflangja{逸失した機会}{lost opportunity}};
  \end{scope}
  % ---------------- (b) elastic capacity
  \begin{scope}[shift={(63,0)}]
    \node[ttl] at (-4,43) {\iflangja{(b) 伸縮する容量}{(b) elastic capacity}};
    \draw[dem] plot[smooth, tension=0.6] coordinates
      {(0,6) (5,8) (10,12) (15,10) (20,14) (25,17) (30,15) (35,19) (38,29) (41,28) (44,18) (50,16)};
    \draw[capa] (0,9) -- (5,9) -- (5,13.5) -- (10,13.5) -- (10,14) -- (15,14)
      -- (15,15.5) -- (20,15.5) -- (20,19) -- (30,19) -- (30,21) -- (35,21)
      -- (35,31) -- (43,31) -- (43,21) -- (46,21) -- (46,18.5) -- (50,18.5);
    \draw[ax] (0,0) -- (53,0) node[below left, lab] {\iflangja{時間}{time}};
    \draw[ax] (0,0) -- (0,35) node[above right, lab, anchor=south west] {\iflangja{資源}{resources}};
  \end{scope}
  % ---------------- legend
  \draw[dem] (36,-7) -- (44,-7);
  \node[lab, anchor=west] at (45,-7) {\iflangja{需要}{demand}};
  \draw[capa] (62,-7) -- (70,-7);
  \node[lab, anchor=west] at (71,-7) {\iflangja{容量}{capacity}};
\end{tikzpicture}
   (width ~118 mm)
\begin{tikzpicture}[x=1mm, y=1mm, font=\scriptsize\sffamily,
  >={Stealth[length=1.5mm]},
  ax/.style={->, thin},
  dem/.style={thick, ln-accent},
  capa/.style={thick, ln-caution},
  lab/.style={font=\scriptsize\sffamily, align=center},
  ttl/.style={font=\footnotesize\sffamily\bfseries, anchor=west}]
  % ---------------- (a) fixed capacity
  \begin{scope}[shift={(3,0)}]
    \node[ttl] at (-4,43) {\iflangja{(a) 固定された容量}{(a) fixed capacity}};
    \begin{scope}
      \clip (0,0) rectangle (50,22);
      \fill[ln-rule!45] (0,34) -- plot[smooth, tension=0.6] coordinates
        {(0,6) (5,8) (10,12) (15,10) (20,14) (25,17) (30,15) (35,19) (38,29) (41,28) (44,18) (50,16)}
        -- (50,34) -- cycle;
    \end{scope}
    \begin{scope}
      \clip (0,22) rectangle (50,34);
      \fill[ln-caution!30] (0,0) -- plot[smooth, tension=0.6] coordinates
        {(0,6) (5,8) (10,12) (15,10) (20,14) (25,17) (30,15) (35,19) (38,29) (41,28) (44,18) (50,16)}
        -- (50,0) -- cycle;
    \end{scope}
    \draw[dem] plot[smooth, tension=0.6] coordinates
      {(0,6) (5,8) (10,12) (15,10) (20,14) (25,17) (30,15) (35,19) (38,29) (41,28) (44,18) (50,16)};
    \draw[capa] (0,22) -- (50,22);
    \draw[ax] (0,0) -- (53,0) node[below left, lab] {\iflangja{時間}{time}};
    \draw[ax] (0,0) -- (0,35) node[above right, lab, anchor=south west] {\iflangja{資源}{resources}};
    \node[lab, fill=white, inner sep=0.6pt] at (14,18.5) {\iflangja{無駄な資源}{wasted resources}};
    \draw[->, ln-caution] (30,30) -- (37,25.5);
    \node[lab, anchor=east, text=ln-caution] at (30.5,30.5) {\iflangja{逸失した機会}{lost opportunity}};
  \end{scope}
  % ---------------- (b) elastic capacity
  \begin{scope}[shift={(63,0)}]
    \node[ttl] at (-4,43) {\iflangja{(b) 伸縮する容量}{(b) elastic capacity}};
    \draw[dem] plot[smooth, tension=0.6] coordinates
      {(0,6) (5,8) (10,12) (15,10) (20,14) (25,17) (30,15) (35,19) (38,29) (41,28) (44,18) (50,16)};
    \draw[capa] (0,9) -- (5,9) -- (5,13.5) -- (10,13.5) -- (10,14) -- (15,14)
      -- (15,15.5) -- (20,15.5) -- (20,19) -- (30,19) -- (30,21) -- (35,21)
      -- (35,31) -- (43,31) -- (43,21) -- (46,21) -- (46,18.5) -- (50,18.5);
    \draw[ax] (0,0) -- (53,0) node[below left, lab] {\iflangja{時間}{time}};
    \draw[ax] (0,0) -- (0,35) node[above right, lab, anchor=south west] {\iflangja{資源}{resources}};
  \end{scope}
  % ---------------- legend
  \draw[dem] (36,-7) -- (44,-7);
  \node[lab, anchor=west] at (45,-7) {\iflangja{需要}{demand}};
  \draw[capa] (62,-7) -- (70,-7);
  \node[lab, anchor=west] at (71,-7) {\iflangja{容量}{capacity}};
\end{tikzpicture}

  \caption{Capacity and demand over time.  (a) A fixed capacity wastes
    resources while demand is low and loses requests during a surge.  (b) An
    elastic capacity follows the demand.}
  \label{fig:l16-capacity}
\end{figure}

Ideally, the capacity of a service would instead follow its demand
(\cref{fig:l16-capacity}b):
\begin{itemize}
  \item capacity scales elastically with the demand, both up and down;
  \item scaling is instantaneous;
  \item the cost is proportional to the demand, and thus to the revenue
    opportunity;
  \item all of this happens automatically, without special expertise;
  \item the resources can be accessed at any time, from anywhere.
\end{itemize}

\begin{definition}[Elasticity]
  The ability of a system to acquire and release resources quickly and
  automatically as its demand rises and falls, so that its capacity follows
  the demand, is called \term{elasticity}[伸縮性].
\end{definition}

The price of this ideal is that a service no longer owns the resources on
which it runs, and its data reside on machines that others operate.

\begin{example}
  The demand of a service is 20 servers during 16 hours of every day and 100
  servers during the remaining 8 hours, $20 \times 16 + 100 \times 8 = 1120$
  server-hours per day.  An owned server costs $c$ per hour, whether it is
  used or not.  Provisioned for the peak, the service pays for
  $100 \times 24 = 2400$ server-hours, $2400c$ per day, and
  $1 - 1120/2400 \approx \SI{53}{\percent}$ of them are wasted.  Provisioned
  for the average demand of $1120/24 \approx 47$ servers, it still wastes
  $27 \times 16 = 432$ server-hours while the demand is low, and during the
  peak it drops $53/100 = \SI{53}{\percent}$ of the requests.  With elastic
  capacity rented at a \SI{50}{\percent} higher price, $1.5c$ per
  server-hour, the service pays $1120 \times 1.5c = 1680c$ per day,
  \SI{30}{\percent} less than with peak provisioning, and drops no requests.
\end{example}

\section{Cloud computing}
\label{sec:l16-cloud}

Cloud computing approximates the ideal described above.\source{P4L4, cloud
computing overview, why cloud computing works, cloud computing vision, cloud
computing definitions; Armbrust et al.\ \cite{armbrust2009}; Mell and Grance
\cite{mell2011}; Silberschatz et al.\ ch.~1.}  Its essential elements are the
following.
\begin{description}
  \item[Shared resources] A provider offers infrastructure---compute, storage
    and network---as well as software and services, and shares them among
    many customers.
  \item[APIs for access and configuration] Customers request, configure and
    release resources through web-based APIs, through libraries for
    programming languages, or with command-line tools.
  \item[Billing and accounting] Usage is metered and billed.  Many pricing
    models coexist, for example on-demand prices, lower prices for capacity
    reserved for a longer period, and \emph{spot} prices for spare capacity
    that vary with supply and demand, so that the provider in effect runs a
    marketplace.  Resources are offered in discrete sizes, such as small,
    medium and extra-large instances with different numbers of CPUs and
    amounts of memory.
  \item[Management by the provider] The provider operates the datacenters,
    the hardware, and the software that implements the services.
\end{description}

\begin{definition}[Cloud computing]
  The provision of computing resources as services over a network, from
  resources that are shared among many customers, acquired and released on
  demand through APIs, billed according to use and managed by the provider,
  is called \term{cloud computing}[クラウドコンピューティング].
\end{definition}

\begin{note}
  Published definitions agree with these elements.  The NIST definition
  \cite{mell2011} names five essential characteristics: on-demand
  self-service, broad network access, resource pooling, rapid elasticity and
  measured service.  Armbrust et al.\ \cite{armbrust2009} use the term for
  both the applications delivered as services over the Internet and the
  hardware and systems software in the datacenters that provide these
  services; the latter they call a \emph{cloud}.
\end{note}

\subsection{Why cloud computing works}

A provider must own enough resources for the demand of all its customers,
and still charge them less than they would spend on resources of their own.
Two principles make this possible.

\begin{definition}[Law of large numbers]
  By the \term{law of large numbers}[大数の法則], the average of $N$
  independent random quantities with the same distribution approaches their
  expected value as $N$ grows.
\end{definition}

The resource needs of a single customer vary widely over time, but different
customers reach their peaks at different times, and the demand averaged over
many customers is roughly constant.  If the demands of $N$ customers are
independent, each with mean $\mu$ and standard deviation $\sigma$, their sum
has mean $N\mu$ and standard deviation $\sigma\sqrt{N}$, so that its
variation relative to its mean, $\sigma/(\mu\sqrt{N})$, shrinks as $N$
grows.  A provider can therefore provision for the aggregate demand, which is
far smaller than the sum of the individual peaks.

\begin{example}
  A provider serves 100 customers.  Each customer needs 10 servers most of
  the time and 100 servers during its peaks, which occupy \SI{10}{\percent}
  of the time and occur independently of the peaks of other customers.
  Provisioning for every customer's peak would take
  $100 \times 100 = \num{10000}$ servers.  The number $K$ of customers at
  their peak at a given moment, however, has mean $100 \times 0.1 = 10$ and
  standard deviation $\sqrt{100 \times 0.1 \times 0.9} = 3$, and the total
  demand, $1000 + 90K$ servers, averages 1900.  With \num{2710} servers,
  enough for $K \le 19$, three standard deviations above the mean, the
  demand exceeds the capacity with a probability of only about
  \SI{0.2}{\percent}.  The provider needs fewer than \SI{30}{\percent} of the
  servers that its customers would need to own.
\end{example}

The second principle is that of \term{economies of scale}[規模の経済]: the
unit cost of providing a resource falls as the quantity provided grows.  A
provider buys hardware, power and network capacity in bulk, spreads the
fixed costs of a datacenter over many machines, and lets each administrator
manage many machines.

\subsection{Computing as a utility}

The vision of cloud computing is older than its technology.  In 1961, John
McCarthy suggested that computing might one day be organized as a public
utility, as the telephone system is, and that such a computer utility could
become the basis of an important new industry.

\begin{definition}[Utility computing]
  In \term{utility computing}[ユーティリティコンピューティング], computing
  resources are provided in the way electricity or telephone service is: as
  a metered, fungible commodity---any unit is as good as any other---that
  customers consume on demand and pay for according to use.
\end{definition}

In the terminology of Armbrust et al.\ \cite{armbrust2009}, utility
computing is what a cloud sells when it is made available to the general
public on a pay-as-you-go basis.  Current clouds realize the vision only in
part, because computing has not become a fully fungible utility.  The
limitations are these:
\begin{description}
  \item[API lock-in] Every provider has its own APIs, and an application
    written for one provider cannot move to another without changes.
  \item[Hardware dependence] Applications may depend on particular hardware,
    such as specific processors or accelerators, which not every provider or
    instance type offers.
  \item[Latency] The resources reside in a remote datacenter, and every
    interaction with them incurs network latency.
  \item[Privacy and security] Data and computations reside on machines that
    the customer neither owns nor controls and that are shared with other
    customers.
\end{description}

\begin{keypoint}
  Cloud computing offers shared, provider-managed resources through APIs,
  elastically and billed by use.  It works because the aggregate demand of
  many customers varies far less than the demand of each of them (law of
  large numbers) and because unit costs fall with scale (economies of
  scale).  It realizes the vision of computing as a utility only in part.
\end{keypoint}

\section{Deployment and service models}
\label{sec:l16-models}

Clouds differ in who may use them and in what they provide.\source{P4L4,
cloud deployment models, cloud service models; Mell and Grance
\cite{mell2011}; Silberschatz et al.\ ch.~1.}  The \emph{deployment models}
describe who uses a cloud.
\begin{description}
  \item[Public cloud] A \term{public cloud}[パブリッククラウド] is operated by
    a provider for third-party customers, its \emph{tenants}, who share its
    resources.
  \item[Private cloud] In a \term{private cloud}[プライベートクラウド] an
    organization applies cloud technology, such as virtualization, APIs and
    automated provisioning, to its own infrastructure, for its own
    applications.
  \item[Hybrid cloud] A \term{hybrid cloud}[ハイブリッドクラウド] combines a
    private cloud with a public one.  An organization runs its applications
    in its private cloud and uses the public cloud for failover, to absorb
    spikes in demand that exceed its own capacity, or for testing.
  \item[Community cloud] A \term{community cloud}[コミュニティクラウド] is a
    public cloud whose use is restricted to a certain type of users, such as
    government agencies or research institutions.
\end{description}

The \emph{service models} describe how much of the stack of hardware and
software the provider manages (\cref{fig:l16-service-models}).  In an
on-premises installation the organization manages the entire stack itself,
from the network, storage and servers through virtualization, the operating
system, middleware and runtime to its data and applications.

\begin{figure}[htbp]
  \centering
  % Cloud service models: which layers of the software and hardware stack the
% customer manages and which the provider manages, for an on-premises
% installation, IaaS, PaaS and SaaS.
% Usage: % Cloud service models: which layers of the software and hardware stack the
% customer manages and which the provider manages, for an on-premises
% installation, IaaS, PaaS and SaaS.
% Usage: % Cloud service models: which layers of the software and hardware stack the
% customer manages and which the provider manages, for an on-premises
% installation, IaaS, PaaS and SaaS.
% Usage: \input{figures/l16-service-models}   (width ~116 mm)
\begin{tikzpicture}[x=1mm, y=1mm, font=\scriptsize\sffamily,
  own/.style={draw, fill=white, minimum width=26mm, minimum height=4.6mm, inner sep=0.5pt},
  prov/.style={draw, fill=ln-accent!25, minimum width=26mm, minimum height=4.6mm, inner sep=0.5pt},
  hd/.style={font=\footnotesize\sffamily\bfseries, anchor=base},
  ex/.style={font=\scriptsize\sffamily, text=ln-muted, anchor=south},
  lab/.style={font=\scriptsize\sffamily, anchor=west}]
  % layers from top (1) to bottom (9)
  \def\layername#1{%
    \ifcase#1\or \iflangja{アプリケーション}{applications}%
    \or \iflangja{データ}{data}%
    \or \iflangja{ランタイム}{runtime}%
    \or \iflangja{ミドルウェア}{middleware}%
    \or OS%
    \or \iflangja{仮想化}{virtualization}%
    \or \iflangja{サーバ}{servers}%
    \or \iflangja{ストレージ}{storage}%
    \or \iflangja{ネットワーク}{networking}\fi}
  % column / header / example / first layer managed by the provider (10 = none)
  \foreach \c/\h/\e/\p in {%
      0/{\iflangja{オンプレミス}{on premises}}/{}/10,
      1/IaaS/{Amazon EC2}/6,
      2/PaaS/{Google App Engine}/3,
      3/SaaS/{Gmail}/1} {
    \pgfmathsetmacro{\cx}{14 + 30*\c}
    \node[hd] at (\cx,52.3) {\h};
    \node[ex] at (\cx,46.8) {\strut\e};
    \foreach \l in {1,...,9} {
      \pgfmathsetmacro{\ly}{46 - 4.6*\l + 2.3}
      \ifnum\l<\p
        \node[own] at (\cx,\ly) {\layername{\l}};
      \else
        \node[prov] at (\cx,\ly) {\layername{\l}};
      \fi
    }
  }
  % legend
  \node[own, minimum width=5mm, minimum height=3.5mm] at (27,-3.5) {};
  \node[lab] at (30,-3.5) {\iflangja{利用者が管理}{managed by the customer}};
  \node[prov, minimum width=5mm, minimum height=3.5mm] at (70,-3.5) {};
  \node[lab] at (73,-3.5) {\iflangja{プロバイダが管理}{managed by the provider}};
\end{tikzpicture}
   (width ~116 mm)
\begin{tikzpicture}[x=1mm, y=1mm, font=\scriptsize\sffamily,
  own/.style={draw, fill=white, minimum width=26mm, minimum height=4.6mm, inner sep=0.5pt},
  prov/.style={draw, fill=ln-accent!25, minimum width=26mm, minimum height=4.6mm, inner sep=0.5pt},
  hd/.style={font=\footnotesize\sffamily\bfseries, anchor=base},
  ex/.style={font=\scriptsize\sffamily, text=ln-muted, anchor=south},
  lab/.style={font=\scriptsize\sffamily, anchor=west}]
  % layers from top (1) to bottom (9)
  \def\layername#1{%
    \ifcase#1\or \iflangja{アプリケーション}{applications}%
    \or \iflangja{データ}{data}%
    \or \iflangja{ランタイム}{runtime}%
    \or \iflangja{ミドルウェア}{middleware}%
    \or OS%
    \or \iflangja{仮想化}{virtualization}%
    \or \iflangja{サーバ}{servers}%
    \or \iflangja{ストレージ}{storage}%
    \or \iflangja{ネットワーク}{networking}\fi}
  % column / header / example / first layer managed by the provider (10 = none)
  \foreach \c/\h/\e/\p in {%
      0/{\iflangja{オンプレミス}{on premises}}/{}/10,
      1/IaaS/{Amazon EC2}/6,
      2/PaaS/{Google App Engine}/3,
      3/SaaS/{Gmail}/1} {
    \pgfmathsetmacro{\cx}{14 + 30*\c}
    \node[hd] at (\cx,52.3) {\h};
    \node[ex] at (\cx,46.8) {\strut\e};
    \foreach \l in {1,...,9} {
      \pgfmathsetmacro{\ly}{46 - 4.6*\l + 2.3}
      \ifnum\l<\p
        \node[own] at (\cx,\ly) {\layername{\l}};
      \else
        \node[prov] at (\cx,\ly) {\layername{\l}};
      \fi
    }
  }
  % legend
  \node[own, minimum width=5mm, minimum height=3.5mm] at (27,-3.5) {};
  \node[lab] at (30,-3.5) {\iflangja{利用者が管理}{managed by the customer}};
  \node[prov, minimum width=5mm, minimum height=3.5mm] at (70,-3.5) {};
  \node[lab] at (73,-3.5) {\iflangja{プロバイダが管理}{managed by the provider}};
\end{tikzpicture}
   (width ~116 mm)
\begin{tikzpicture}[x=1mm, y=1mm, font=\scriptsize\sffamily,
  own/.style={draw, fill=white, minimum width=26mm, minimum height=4.6mm, inner sep=0.5pt},
  prov/.style={draw, fill=ln-accent!25, minimum width=26mm, minimum height=4.6mm, inner sep=0.5pt},
  hd/.style={font=\footnotesize\sffamily\bfseries, anchor=base},
  ex/.style={font=\scriptsize\sffamily, text=ln-muted, anchor=south},
  lab/.style={font=\scriptsize\sffamily, anchor=west}]
  % layers from top (1) to bottom (9)
  \def\layername#1{%
    \ifcase#1\or \iflangja{アプリケーション}{applications}%
    \or \iflangja{データ}{data}%
    \or \iflangja{ランタイム}{runtime}%
    \or \iflangja{ミドルウェア}{middleware}%
    \or OS%
    \or \iflangja{仮想化}{virtualization}%
    \or \iflangja{サーバ}{servers}%
    \or \iflangja{ストレージ}{storage}%
    \or \iflangja{ネットワーク}{networking}\fi}
  % column / header / example / first layer managed by the provider (10 = none)
  \foreach \c/\h/\e/\p in {%
      0/{\iflangja{オンプレミス}{on premises}}/{}/10,
      1/IaaS/{Amazon EC2}/6,
      2/PaaS/{Google App Engine}/3,
      3/SaaS/{Gmail}/1} {
    \pgfmathsetmacro{\cx}{14 + 30*\c}
    \node[hd] at (\cx,52.3) {\h};
    \node[ex] at (\cx,46.8) {\strut\e};
    \foreach \l in {1,...,9} {
      \pgfmathsetmacro{\ly}{46 - 4.6*\l + 2.3}
      \ifnum\l<\p
        \node[own] at (\cx,\ly) {\layername{\l}};
      \else
        \node[prov] at (\cx,\ly) {\layername{\l}};
      \fi
    }
  }
  % legend
  \node[own, minimum width=5mm, minimum height=3.5mm] at (27,-3.5) {};
  \node[lab] at (30,-3.5) {\iflangja{利用者が管理}{managed by the customer}};
  \node[prov, minimum width=5mm, minimum height=3.5mm] at (70,-3.5) {};
  \node[lab] at (73,-3.5) {\iflangja{プロバイダが管理}{managed by the provider}};
\end{tikzpicture}

  \caption{Cloud service models, with an example of each.  The shaded layers
    are managed by the provider, the others by the customer.}
  \label{fig:l16-service-models}
\end{figure}

With
\term{infrastructure as a service}[IaaS]\index{IaaS|see{infrastructure as a service}}
(IaaS) the provider manages the physical infrastructure, that is, the
network, storage and servers, and its virtualization.  Customers rent virtual machines and
install and manage the operating system and everything above it.  Amazon
EC2 is an example.

With
\term{platform as a service}[PaaS]\index{PaaS|see{platform as a service}}
(PaaS) the provider also manages the operating system, the middleware and
the runtime environment.  Customers supply only their applications, written
against the interfaces of the platform, and their data, and the platform
deploys and runs the applications.  Google App Engine is an example.

With
\term{software as a service}[SaaS]\index{SaaS|see{software as a service}}
(SaaS) the provider manages the entire stack, the application included, and
customers merely use the application, as the users of Gmail do.

\begin{example}
  The operator of the three-tier online shop of \cref{sec:l16-services} can
  run it on IaaS by renting virtual machines, installing an operating
  system, a web server, an application server and a database server on them,
  and managing all this software, its updates and its scaling.  On PaaS the
  operator uploads only the code of the business logic and the static
  content; the platform runs them, provides a managed database, and adds
  instances as the load grows.  The employees of the shop may at the same
  time read their e-mail through a SaaS offering, for which they manage
  nothing but their accounts and messages.
\end{example}

\section{Requirements for the cloud}
\label{sec:l16-requirements}

From the provider's point of view, cloud computing places a number of
requirements on the underlying technology.\source{P4L4, requirements for the
cloud, cloud failure probability; OSTEP ch.~48.}
\begin{enumerate}
  \item \textbf{Fungible resources.}  Resources must be interchangeable and
    easy to repurpose for customers with different requirements, such as
    different operating systems and software stacks.  Without fungibility
    every customer would need dedicated resources, and the economic
    opportunity of the cloud, which rests on sharing
    (\cref{sec:l16-cloud}), would not exist.
  \item \textbf{Elastic, dynamic resource allocation.}  Mechanisms and
    policies must allocate resources on demand and reclaim them when they
    are released.
  \item \textbf{Scale.}  Management and resource allocation must work for
    datacenters of many thousands of machines and for many customers.
  \item \textbf{Dealing with failures.}  At this scale failures are the
    normal case, and the system must continue to operate despite them.
  \item \textbf{Multi-tenancy.}  Tenants must be isolated from each other,
    and the performance of each of them must be guaranteed.
  \item \textbf{Security.}  The data and computations of every tenant must be
    protected from other tenants and from attackers outside the cloud.
\end{enumerate}

\begin{definition}[Multi-tenancy]
  The sharing of the same physical resources by several customers, each of
  which uses its own, isolated set of resources, is called
  \term{multi-tenancy}[マルチテナンシ].
\end{definition}

Multi-tenancy requires isolation\index{isolation} of the tenants' state, as
a VMM provides it for its VMs (Lesson~12), but also isolation of their
performance.

\begin{definition}[Performance isolation]
  A system provides \term{performance isolation}[性能隔離] when the
  performance that a tenant obtains does not depend on the workloads of the
  other tenants that share its resources.
\end{definition}

Without performance isolation, a tenant whose VM shares a machine with a VM
that saturates the disk or the memory observes unpredictable response times,
and the provider cannot keep the service level
agreements\index{service level agreement} (Lesson~6) it has promised.

\subsection{Failures at scale}

Failures become frequent simply because of the number of components.  If
each of $N$ components fails independently with probability $p$ within a
given period, the probability that at least one of them fails is
\begin{equation}
  P_{\mathrm{fail}} \;=\; 1 - (1-p)^{N} ,
  \label{eq:l16-failure}
\end{equation}
which approaches 1 as $N$ grows, while the expected number of failed
components, $Np$, grows linearly with~$N$.  A component that rarely fails
on its own therefore fails somewhere in a datacenter all the time.  Systems
in the cloud must detect failures, for instance with timeouts, and continue
despite them, with the mechanisms of the preceding lessons: replicas that
take over for failed ones (Lesson~14), checkpoints from which computations
resume (Lesson~8)\index{checkpointing}, and migration of VMs away from
failing machines (Lesson~12).

\begin{example}
  A machine fails on a given day with probability $p = 0.001$, independently
  of the other machines; on its own, it fails on average once in about three
  years.  By \cref{eq:l16-failure}, the probability that at least one
  machine fails on a given day is \SI{1.0}{\percent} for 10 machines,
  \SI{9.5}{\percent} for 100 machines and \SI{63}{\percent} for \num{1000}
  machines.  For \num{10000} machines it is practically \SI{100}{\percent},
  and on average 10 machines fail every day.
\end{example}

\begin{keypoint}
  A cloud requires fungible resources, elastic allocation, management at
  scale, tolerance of failures, and isolation of the state and the
  performance of its tenants.  With many components, some component is
  almost always failing (\cref{eq:l16-failure}).
\end{keypoint}

\section{Cloud enabling technologies}
\label{sec:l16-technologies}

Several technologies together make cloud computing possible.\source{P4L4,
cloud enabling technologies; Hindman et al.\ \cite{hindman2011}; Dean and
Ghemawat \cite{dean2004}; Zaharia et al.\ \cite{zaharia2012}; McKeown et
al.\ \cite{mckeown2008}.}

\paragraph{Virtualization.}  Virtualization (Lesson~12) makes resources
fungible.  A physical machine can host the VMs of any tenants, with any
operating systems and software stacks, and the VMs isolate the tenants from
each other.  VMs can be created, destroyed and migrated as the demand
changes.

\paragraph{Resource provisioning.}  Resource managers for clusters, such as
Mesos \cite{hindman2011} and Hadoop YARN, schedule the resources of many
machines among the applications and frameworks that run on them, deciding
which CPUs and how much memory on which machines each of them receives.

\paragraph{Big data processing.}  Frameworks such as Hadoop MapReduce and
Spark process large data sets in parallel on many machines.

\begin{definition}[MapReduce]
  \term{MapReduce}[MapReduce] \cite{dean2004} is a programming model, and a
  runtime system that implements it, for processing large data sets on a
  cluster.  The programmer writes a \emph{map} function, which transforms
  each input record into intermediate key/value pairs, and a \emph{reduce}
  function, which combines all intermediate values with the same key.  The
  runtime splits the input, runs map and reduce tasks in parallel on many
  machines, groups the intermediate pairs by key, and re-executes the tasks
  of machines that fail.
\end{definition}

To count how often every word occurs in a collection of documents, for
instance, the map function emits the pair $(w, 1)$ for every occurrence of a
word~$w$, and the reduce function sums the values for each word.  MapReduce
stores intermediate results on disk; Spark \cite{zaharia2012} keeps them in
memory, which speeds up iterative algorithms and interactive queries that
reuse the same data.

\paragraph{Storage.}  Clouds store data in distributed file systems
(Lesson~14) designed for very large files that are mostly appended to rather
than overwritten; in NoSQL databases, which scale out by giving up the
relational data model and some of the consistency guarantees of traditional
databases; and in distributed in-memory stores and caches, which keep data
in the memory of many machines.

\paragraph{Software-defined networking, storage and datacenters.}
Networks, storage and eventually entire datacenters are configured by
software through APIs rather than by hand.

\begin{definition}[Software-defined networking]
  In
  \term{software-defined networking}[SDN]\index{SDN|see{software-defined networking}}
  (SDN), the \emph{control plane} of a network, which decides how traffic is
  forwarded, is separated from the \emph{data plane} of the switches, which
  forwards the packets.  A logically centralized controller, implemented in
  software, programs the forwarding tables of the switches through an open
  interface such as OpenFlow \cite{mckeown2008}.
\end{definition}

With SDN a cloud can reconfigure its network by software, for example to
give every tenant a virtual network of its own that is isolated from those
of other tenants.  Software-defined storage and software-defined datacenters
apply the same principle to storage and to all resources of a datacenter.

\paragraph{Monitoring.}  Thousands of machines and the services running on
them continuously produce logs and measurements.  Real-time log processing
systems, such as Flume, CloudWatch or Log Insight, collect and analyze them
as they are produced, in order to detect failures, anomalies and overloads,
and to trigger elastic scaling.

\section{The cloud as a big data engine}
\label{sec:l16-bigdata}

Many applications analyze data sets that no single machine can hold, such as
the logs of user activity of a large Internet service, web pages to be
indexed, or scientific measurements.\source{P4L4, the cloud as a big data
engine, example big data stacks.}

\begin{definition}[Big data]
  Data sets whose size exceeds what a single machine can store and process,
  so that they must be stored and analyzed on many machines, are called
  \term{big data}[ビッグデータ].
\end{definition}

The cloud provides both the storage for such data and the many machines that
process them in parallel.  The software that does so forms a \emph{big data
engine} with the following layers (\cref{fig:l16-bigdata}).
\begin{description}
  \item[Data storage] stores the data, distributed and replicated over many
    machines.
  \item[Data processing] executes computations over the data in parallel on
    many machines.
  \item[Caching] keeps frequently used data in memory, avoiding the latency of
    the storage layer.
  \item[Language front-ends] let users query and analyze the data in
    high-level languages, for instance SQL-like query languages, instead of
    programming the processing layer directly.
  \item[Analytics libraries] provide algorithms built on the processing
    layer, for instance for machine learning, graph analysis or statistics.
  \item[Streaming data] such as logs, clicks or sensor readings, enter the
    engine continuously and must be processed as they arrive, not only after
    they have been stored.
\end{description}

\begin{figure}[htbp]
  \centering
  % The cloud as a big data engine: the layers of a big data stack, and the
% streaming data that enter it continuously.
% Usage: % The cloud as a big data engine: the layers of a big data stack, and the
% streaming data that enter it continuously.
% Usage: % The cloud as a big data engine: the layers of a big data stack, and the
% streaming data that enter it continuously.
% Usage: \input{figures/l16-bigdata}   (width ~116 mm)
\begin{tikzpicture}[x=1mm, y=1mm, font=\scriptsize\sffamily,
  >={Stealth[length=1.5mm]},
  ly/.style={draw, fill=ln-box, minimum height=8mm, inner sep=1pt, align=center},
  top/.style={draw, fill=ln-accent!18, minimum height=8mm, inner sep=1pt, align=center},
  st/.style={draw, fill=ln-accent!35, minimum height=8mm, inner sep=1pt, align=center},
  lab/.style={font=\scriptsize\sffamily, text=ln-muted, align=center}]
  % users of the engine
  \node[lab] at (74,49) {\iflangja{アプリケーション，データ分析者}{applications, data analysts}};
  \draw[<->, thick] (57,46.5) -- (57,42.5);
  \draw[<->, thick] (95,46.5) -- (95,42.5);
  % layers
  \node[top, minimum width=42mm] (lang) at (53,38) {\iflangja{言語フロントエンド}{language front-ends}\\\textcolor{ln-muted}{\iflangja{例: 問い合わせ}{e.g.\ queries}}};
  \node[top, minimum width=42mm] (lib) at (97,38) {\iflangja{分析ライブラリ}{analytics libraries}\\\textcolor{ln-muted}{\iflangja{例: 機械学習}{e.g.\ machine learning}}};
  \node[ly, minimum width=86mm] (proc) at (75,27) {\iflangja{データ処理}{data processing}\\\textcolor{ln-muted}{\iflangja{多数のマシンで並列に実行}{in parallel on many machines}}};
  \node[ly, minimum width=86mm] (cache) at (75,16) {\iflangja{キャッシュ}{caching}\\\textcolor{ln-muted}{\iflangja{頻繁に使うデータをメモリに保持}{frequently used data kept in memory}}};
  \node[st, minimum width=86mm] (store) at (75,5) {\iflangja{データストレージ}{data storage}\\\textcolor{ln-muted}{\iflangja{多数のマシンに分散して格納}{distributed over many machines}}};
  % streaming data
  \node[lab, text=ln-accent, anchor=east] (sd) at (27,27)
    {\iflangja{連続して到着する}{continuously}\\\iflangja{ストリームデータ}{streaming data}};
  \draw[->, line width=1.2pt, ln-accent] (sd.east) -- (proc.west);
  \draw[->, line width=1.2pt, ln-accent] (29.2,27) |- (store.west);
  \fill[ln-accent] (29.2,27) circle (0.8);
\end{tikzpicture}
   (width ~116 mm)
\begin{tikzpicture}[x=1mm, y=1mm, font=\scriptsize\sffamily,
  >={Stealth[length=1.5mm]},
  ly/.style={draw, fill=ln-box, minimum height=8mm, inner sep=1pt, align=center},
  top/.style={draw, fill=ln-accent!18, minimum height=8mm, inner sep=1pt, align=center},
  st/.style={draw, fill=ln-accent!35, minimum height=8mm, inner sep=1pt, align=center},
  lab/.style={font=\scriptsize\sffamily, text=ln-muted, align=center}]
  % users of the engine
  \node[lab] at (74,49) {\iflangja{アプリケーション，データ分析者}{applications, data analysts}};
  \draw[<->, thick] (57,46.5) -- (57,42.5);
  \draw[<->, thick] (95,46.5) -- (95,42.5);
  % layers
  \node[top, minimum width=42mm] (lang) at (53,38) {\iflangja{言語フロントエンド}{language front-ends}\\\textcolor{ln-muted}{\iflangja{例: 問い合わせ}{e.g.\ queries}}};
  \node[top, minimum width=42mm] (lib) at (97,38) {\iflangja{分析ライブラリ}{analytics libraries}\\\textcolor{ln-muted}{\iflangja{例: 機械学習}{e.g.\ machine learning}}};
  \node[ly, minimum width=86mm] (proc) at (75,27) {\iflangja{データ処理}{data processing}\\\textcolor{ln-muted}{\iflangja{多数のマシンで並列に実行}{in parallel on many machines}}};
  \node[ly, minimum width=86mm] (cache) at (75,16) {\iflangja{キャッシュ}{caching}\\\textcolor{ln-muted}{\iflangja{頻繁に使うデータをメモリに保持}{frequently used data kept in memory}}};
  \node[st, minimum width=86mm] (store) at (75,5) {\iflangja{データストレージ}{data storage}\\\textcolor{ln-muted}{\iflangja{多数のマシンに分散して格納}{distributed over many machines}}};
  % streaming data
  \node[lab, text=ln-accent, anchor=east] (sd) at (27,27)
    {\iflangja{連続して到着する}{continuously}\\\iflangja{ストリームデータ}{streaming data}};
  \draw[->, line width=1.2pt, ln-accent] (sd.east) -- (proc.west);
  \draw[->, line width=1.2pt, ln-accent] (29.2,27) |- (store.west);
  \fill[ln-accent] (29.2,27) circle (0.8);
\end{tikzpicture}
   (width ~116 mm)
\begin{tikzpicture}[x=1mm, y=1mm, font=\scriptsize\sffamily,
  >={Stealth[length=1.5mm]},
  ly/.style={draw, fill=ln-box, minimum height=8mm, inner sep=1pt, align=center},
  top/.style={draw, fill=ln-accent!18, minimum height=8mm, inner sep=1pt, align=center},
  st/.style={draw, fill=ln-accent!35, minimum height=8mm, inner sep=1pt, align=center},
  lab/.style={font=\scriptsize\sffamily, text=ln-muted, align=center}]
  % users of the engine
  \node[lab] at (74,49) {\iflangja{アプリケーション，データ分析者}{applications, data analysts}};
  \draw[<->, thick] (57,46.5) -- (57,42.5);
  \draw[<->, thick] (95,46.5) -- (95,42.5);
  % layers
  \node[top, minimum width=42mm] (lang) at (53,38) {\iflangja{言語フロントエンド}{language front-ends}\\\textcolor{ln-muted}{\iflangja{例: 問い合わせ}{e.g.\ queries}}};
  \node[top, minimum width=42mm] (lib) at (97,38) {\iflangja{分析ライブラリ}{analytics libraries}\\\textcolor{ln-muted}{\iflangja{例: 機械学習}{e.g.\ machine learning}}};
  \node[ly, minimum width=86mm] (proc) at (75,27) {\iflangja{データ処理}{data processing}\\\textcolor{ln-muted}{\iflangja{多数のマシンで並列に実行}{in parallel on many machines}}};
  \node[ly, minimum width=86mm] (cache) at (75,16) {\iflangja{キャッシュ}{caching}\\\textcolor{ln-muted}{\iflangja{頻繁に使うデータをメモリに保持}{frequently used data kept in memory}}};
  \node[st, minimum width=86mm] (store) at (75,5) {\iflangja{データストレージ}{data storage}\\\textcolor{ln-muted}{\iflangja{多数のマシンに分散して格納}{distributed over many machines}}};
  % streaming data
  \node[lab, text=ln-accent, anchor=east] (sd) at (27,27)
    {\iflangja{連続して到着する}{continuously}\\\iflangja{ストリームデータ}{streaming data}};
  \draw[->, line width=1.2pt, ln-accent] (sd.east) -- (proc.west);
  \draw[->, line width=1.2pt, ln-accent] (29.2,27) |- (store.west);
  \fill[ln-accent] (29.2,27) circle (0.8);
\end{tikzpicture}

  \caption{The layers of a big data engine.  Streaming data enter the
    engine continuously; they are processed as they arrive and stored.}
  \label{fig:l16-bigdata}
\end{figure}

Two widely used open-source stacks implement these layers
(\cref{tab:l16-stacks}).  The \emph{Hadoop} stack is built around the
Hadoop Distributed File System (HDFS) and MapReduce.  The \emph{Berkeley
Data Analytics Stack} (BDAS), developed at the University of California,
Berkeley, is built around Spark and the cluster resource manager Mesos, and
it can use HDFS and other storage systems below it.

\begin{table}[htbp]
  \centering
  \caption{Components of two big data stacks, by layer.}
  \label{tab:l16-stacks}
  \small
  \begin{tabular}{@{}>{\raggedright\arraybackslash}p{0.2\linewidth}>{\raggedright\arraybackslash}p{0.36\linewidth}>{\raggedright\arraybackslash}p{0.36\linewidth}@{}}
    \toprule
    Layer & Hadoop & BDAS \\
    \midrule
    Data storage & HDFS (file system); HBase (column-oriented store) & HDFS,
      Amazon S3 and other storage systems \\
    Caching & --- & Tachyon (in-memory storage) \\
    Resource management & YARN & Mesos, YARN \\
    Data processing & MapReduce & Spark \\
    Streaming data & Flume (log collection) & Spark Streaming \\
    Language front-ends & Hive (SQL), Pig (data flow) & Shark (SQL), BlinkDB
      (approximate queries) \\
    Analytics libraries & Mahout (machine learning), R connectors
      (statistics) & MLlib, MLbase (machine learning), GraphX (graphs) \\
    Coordination & ZooKeeper (coordination), Oozie (workflows), Sqoop (data
      exchange with relational databases) & --- \\
    \bottomrule
  \end{tabular}
\end{table}

\needspace{8\baselineskip}
\begin{keypoint}[Lesson summary]
  Internet services consist of presentation, business logic and database
  tiers, and scale out over many nodes behind a load balancer; homogeneous
  nodes are simple to manage, while heterogeneous nodes exploit locality but
  need request-aware front ends and demand-driven management.  Cloud
  computing provides shared resources elastically through APIs, billed by
  use and managed by the provider.  It works because of the law of large
  numbers and economies of scale, and it approaches the vision of utility
  computing.  Clouds are public, private, hybrid or community clouds, and
  offer infrastructure, platforms or software as a service.  They require
  fungible resources, elastic allocation at scale, tolerance of frequent
  failures and isolation of their tenants, which virtualization, cluster
  resource management, big data frameworks such as MapReduce, distributed
  storage, software-defined networking and monitoring provide.  On these
  technologies the cloud serves as a big data engine.
\end{keypoint}

\end{document}
